\ssscp{Schwa}{03-02-01-02}
As discussed in \srf{01-07},
\tsc{PAn} is reconstructed as having a four-vowel system with
*i, *ə, *a, *u \citep[173]{Blust2013}.
Most Austronesian languages in other Wallacean subgroups tend
to have a basic five-vowel system with /i/,
/e/, /a/, /o/, /u/.
However most languages of Flores have six-vowel systems
with these five vowels and the addition of schwa /ə/.
This occurs in Manggarai \citep{VerheijenGrimes1995},
Ngadha \citep{DjawanaiGrimes1995}, Lio \citep{Elias2018},
Palu{\Q}e, and Sika \citep{LewisGrimes1995}.\footnote{
		Havu and Dhao also have schwa which is discussed separately in \srf{03-03-02}.}
Schwa often has several properties in languages of Flores
which set it apart from other vowels in these languages.
Common properties of schwa in these languages are listed below:

\begin{itemize}
	\item The articulation of phonetic schwa is often of short duration and light syllable weight.
	\item Schwa is often restricted to the penultimate syllable.
	\item Schwa does not occur in vowel sequences.
	\item While penultimate stress is the normal rule for Wallacean languages,
				words with penultimate schwa can have stress on other syllables.
	\item In words with an intervocalic consonant /əCV/,
				the consonant following the penultimate schwa
				/ə/ often undergoes phonetic lengthening [əCːV].
\end{itemize}

Some other \tsc{Bima-Lembata} languages also have six-vowel systems,
though lack a true schwa.
Thus, Kedang has lost surface/phonetic schwa but has developed
a contrastive low front vowel /æ/ in addition to back /a/.
Many, but not all, instances of low-front /æ/ are reflexes of
earlier schwa \citep[211ff]{Fricke2019}.

In the languages on Flores,
schwa /ə/ tends to have very light syllabic weight,
be very short in duration,
and tends not to be stress bearing.
So phonemic /əma/ ‘father’ may be realised phonetically as [ə̆ˈmːa]
and may be written orthographically by some speakers as ‹mma›.
The articulation of the initial schwa may be almost inaudible.

%In languages like Ngadha “schwa /ə/ never occurs in vowel sequences”
%\citep[595]{DjawanaiGrimes1995}.
%But in Havu and Dhao, schwa frequently occurs in vowel sequences
%with high vowels as the second vowel (See \trf{02-10}),
%though this is still a restricted distribution compared with
%other vowels which can occur as part of a sequence with all other vowels.
%
%In Ngadha schwa is described as: “The schwa /ə/ is of short duration
%and phonetically lengthens the following consonant.
%Schwa never occurs in the ultimate syllable of disyllabic or trisyllabic lexical roots.
%Root-initial schwa is deleted leaving only the trace lengthening
%of the following consonant.” \citep[596]{DjawanaiGrimes1995}.
%In Palu{\Q}e schwa is described by Donohue (in \citep{GreenhillGray2008}
%note for Palu{\Q}e [Nitung]) as: “ə is an epenthetic,
%unstressed vowel that breaks up CC sequences.”
%
%In Manggarai schwa is described as follows: “Schwa occurs in
%both the ultimate and penultimate syllables of roots that have a medial consonant.
%Schwa does not occur in vowel sequences.”
%and “words are stressed on the penultimate syllable of the root.
%If, however, the penultimate syllable contains a schwa /ə/ that
%is not followed by a prenasalised stop,
%then the ultimate syllable is stressed.” \citep[587f]{VerheijenGrimes1995}.
%
%In Sika schwa is described as: “Syllables with schwa /ə/ as their
%peak are of shorter duration than syllables with /e/ as their peak.
%Schwa occurs in ultimate as well as penultimate syllables.
%[{\ldots}] The vowel /ə/ occurs only rarely in non-geminate vowel
%sequences.” \citep[604]{LewisGrimes1995}.

In a number of languages on Flores there is a retention of PMP
*ə = \it{ə}, particularly in penultimate syllables.
There may also be some lexically specific (kin) terms in which
other vowels become schwa.
This is illustrated in \trf{02-08b}.

\begin{table}
	\centering\tcp{Origins of schwa in some languages of Flores}{02-08b}
	\st{4pt}\begin{tabular}{llllllll}\lsptoprule
gloss&hear&sand&buy&new&egg& three&father\\
PMP&*d\tbr{ə}ŋ\tbr{ə}R&*q\tbr{ə}nay&*b\tbr{ə}li&*baq\tbr{ə}Ru&*qat\tbr{ə}luR&*t\tbr{ə}lu&*ama\\ \midrule
Manggarai&{\hr}d\tbr{ə}ŋe&\tcg{(laiŋ)}&{\hr}ʋ\tbr{ə}li&{\hr}ʋ\tbr{ə}ru&\hp{*qa}t\tbr{ə}lo&{\hr}t\tbr{ə}lu&{\hr}\tbr{ə}ma\\
Waerana&{\hr}d\tbr{ə}ŋeʔ&{\hr}\tbr{ə}nar&{\hr}w\tbr{ə}li&{\hr}w\tbr{ə}ru&\hp{*qa}t\tbr{ə}loʔ&{\hr}t\tbr{ə}luʔ&{\hr}ame\\
Ngadha&{\hr}z\tbr{ə}ŋe&{\hr}\tbr{ə}na&{\hr}v\tbr{ə}li&\tcg{(muzi)}&\hp{*qa}t\tbr{ə}lo&{\hr}t\tbr{ə}lu&{\hr}\tbr{ə}ma\\
Lio&\tcg{(lele)}&{\hr}\tbr{ə}na&\tcg{(gəti)}&\tcg{(muri)}&\hp{*qa}t\tbr{ə}lo&{\hr}t\tbr{ə}lu&{\hr}\tbr{ə}ma\\
Palu{\Q}e&\tcg{(tee)}&\tcg{(kəri)}&&\tcg{(muri)}&\hp{*qa}t\tbr{ə}lo&{\hr}t\tbr{ə}lu&{\hr}ama\\
Sika&{\hr}r\tbr{ə}na&\tcg{(nee)}&\tcg{(βotər)}&\mbox{{\hr}β\tbr{ə}ru-ŋ,}&\hp{*qa}t\tbr{ə}lo&{\hr}t\tbr{ə}lu&{\hr}ama\\
&&&&\hp{*}βoo&&&\\
Kalikasa&{\hr}d\tbr{ə}ŋ\tbr{ə}r&\tcg{(botan)}&\tcg{(opi)}&{\hr}v\tbr{ə}run&\hp{*qa}t\tbr{ə}luk,&{\hr}t\tbr{ə}lu&{\hr}ama\\
&&&&&\hp{*qa}teluk&&\\
		\lspbottomrule
	\end{tabular}
\end{table}
